\section{Award-Layer Triage}
\label{sec:award_screen}

The award layer observes outcomes before it observes conduct. A
losing bid may be sincere, poorly priced, capacity constrained,
exploratory, strategic, or coordinated. The screen therefore
uses what the award layer can measure at the triage stage: which
firms appeared, which firms won, and which firms kept appearing
without ever winning. The object is a ranking for forensic
priority, not proof of conduct or a classification of cartel
membership.

This section defines that ranking and the empirical implications
that discipline it. Subsection~\ref{sec:award_primitive}
introduces persistent zero-win participation as the award-layer
primitive. Subsection~\ref{sec:award_rule} explains why the
continuous ranking can be converted into an operational queue or
binary frequent-loser flag. Subsection
\ref{sec:testable_implications} states the tests that the rest
of the paper uses to evaluate whether the ranking is informative
for enforcement.

\subsection{Persistent Zero-Win Participation}
\label{sec:award_primitive}

Let $T_i$ denote the number of BEC tender-items in which firm
$i$ participates, and let $W_i$ denote the number it wins. In
Preg\~ao, a firm may revise its offer repeatedly within one
tender-item; $T_i$ counts tender-item participations, so these
iterative offers do not raise $T_i$.% CO-AUTHOR EDIT: clarify T_i counts tender-item participations, not Pregao offer revisions (consistent with sec02.2 footnote)
The award-layer screen operates inside the zero-win stratum,
$\{i:W_i=0\}$, which contains $\valAlwaysLosers$ firms. Within
that stratum, the continuous score is
\[
  s_i=\log(1+T_i).
\]
The log transformation stabilizes the scale but does not change
the ordering. The economic object is the rank: among firms the
award layer observes only as losers, which ones appear
persistently enough to warrant bid-layer follow-up?

The maintained condition is monotonic rather than classificatory.
If cartel-deployed losing participants are present, they should
appear more persistently in the zero-win stratum than ordinary
losing firms, because such a role is useful only through repeated
deployment in cartel-targeted tenders. The condition is
maintained, not derived. \ref{app:framework_submission}
formalizes it as a monotone-likelihood-ratio ranking result. The
main text uses it more modestly: persistent zero-win participation
is a cheap statistic that can rank loser-side exposure under a
transparent behavioral assumption and direct attention to the bid
layer where conduct can be evaluated. It does not assign a firm
type or cartel membership.

The zero-win restriction is therefore a scope restriction. Many
firms never win for ordinary competitive reasons: poor fit, high costs,
capacity constraints, market exploration, subcontracting
positioning, or simple bad luck. The screen's content comes from
persistence inside that heterogeneous loser-side set. If the
ranking has enforcement value, it should concentrate
adjudication-anchored loser-side adjacency beyond what generic
participation volume, product-market overlap, and opportunity to
meet CADE anchors would mechanically produce.

\subsection{From Ranking to an Operational Queue}
\label{sec:award_rule}

The continuous score is the central object because agencies can
translate a rank into different investigative protocols. A team
with fixed capacity could inspect the top $k$ firms. A
monitoring unit could refresh the rank each quarter and examine
bunching below old cutoffs. A case team could use the score to
request LANCES bid files for a survivor pool before applying
bid-distribution forensics. These implementations differ, but
all preserve the same sequencing logic: award records first,
costly bid recovery second.% TODO_JLEO_RR_COST_FRONTIER: the later cost-recall-frontier prompt will evaluate multiple first-stage K1 cutoffs along a cost-recall frontier.

For an auditable binary implementation, we also report a
frequent-loser flag. The rule marks zero-win firms whose
participation count exceeds the median plus $1.5$ times the
interquartile range of the always-loser participation
distribution. In BEC, the statistical cutoff is
$\valThresholdStat$, implemented as $T_i\geq\valThreshold$, which
flags $\valFL$ firms, or $\valFLrateAL$ of the always-loser
stratum. The cutoff is fixed from the participation distribution
alone, before any CADE label enters the analysis.

The $T_i\geq\valThreshold$ rule is an auditable administrative
implementation of the rank, not a structural or legal
threshold.% CO-AUTHOR EDIT: make explicit the T_i>=14 rule is an auditable implementation, not a structural threshold
The paper does not
rest on the number $14$; it rests on the continuous loser-side
participation ranking. Two analysts applying the same rule to
the same award records obtain the same first-stage list. Above
the cutoff does not mean cartel member; below it does not mean
clean. Because a published static threshold can be gamed,
strategic adaptation makes continuous ranks and periodically
refreshed queues preferable to a permanent static cutoff: the
enforcement architecture treats the binary flag as one
implementation among several, alongside continuous ranks,
top-$k$ queues, refreshed thresholds, and budget-calibrated
survivor pools.% CO-AUTHOR EDIT: strategic adaptation favors continuous ranks / refreshed queues over a permanent static cutoff
The strategic-adaptation analysis later returns to this point.

\subsection{Testable Implications}
\label{sec:testable_implications}

The screen is useful only if it survives the threats implied by
its limited information. Table~\ref{tab:testable_implications_submission}
is a forward-looking validation-threat map: for each implication
it states the required test and what would weaken the design, not
just what the design predicts. The implications it lists are the
checks the rest of the paper must clear, not completed evidence.
The first and core implication is the validation claim:
% TODO_JLEO_RR_OPPORTUNITY_TABLE: \S4.2 opportunity-adjusted validation (Table B) is the empirical answer to the mechanical-exposure threat flagged here.
persistent zero-win participation should rank
adjudication-anchored loser-side cobidders within the always-loser
stratum. The main threat is mechanical exposure: frequent
participants, and firms operating in the same procurement
environments as CADE defendants, have more opportunities to
appear near them even under ordinary competition.% CO-AUTHOR EDIT: foreground core implication first; name mechanical exposure (volume + opportunity set) as main threat; call Table 2 a validation-threat map

The second implication is a scope check. A loser-side statistic
should perform worse against direct CADE defendants than against
always-loser cobidders, because it ranks a different role in the
procurement environment. The third implication asks whether the
cobidder target has economic content: among frequent losers,
cobidders should differ from other frequent losers in directions
consistent with loser-side exposure, proximity to defendants,
market concentration, and bid-level losing behavior. These
profiles give the validation target economic content; they do
not prove conduct or assign legal status.

The final implications evaluate the enforcement architecture.
If award records and bid records carry different information,
the award-layer score should add non-redundant signal to
bid-distribution forensics. If the screen has institutional
value, using it before bid recovery should reduce the
bid-microdata footprint while preserving much of the
adjudication-anchored recall. And because firms can respond to
published administrative rules, the binary threshold should be
treated as vulnerable to strategic adaptation rather than as a
fixed rule immune to gaming.

% CO-AUTHOR EDIT: rebuilt Table 2 as a 6-column x 10-row pre-analysis validation-threat map (test names + metric names + qualitative pass/weaken interpretations only; no result values)
\begin{table}[htbp]
\centering
\caption{Validation-threat map: implications, required tests, and interpretation.}
\label{tab:testable_implications_submission}
\scriptsize
\begin{adjustbox}{max width=\textwidth,center}
\begin{threeparttable}
\begin{tabular}{p{1.9cm}p{2.3cm}p{2.4cm}p{2.2cm}p{2.5cm}p{2.5cm}}
\toprule
Threat / implication & Why it matters & Required test & Main metric & Interpretation if test passes & Interpretation if test weakens \\
\midrule
Label construction & Validation target is the paper's ground truth. & Label funnel and count reconciliation. & Cases, defendants, and cobidders by benchmark. & Target counts trace to data. & Sample-credibility problem. \\
Participation volume & Screen could be high-volume losing only. & Volume controls and participation-matched permutations. & PR-AUC, lift, AUC, excess ratio. & Signal beyond participation volume. & Screen is volume only. \\
Opportunity-set exposure & Firms in the same item, buyer, year, and modality cells as defendants meet them more. & Expected-contact within opportunity cells. & PR-AUC, precision@$k$, recall@$k$, excess contact. & Signal beyond exposure. & Main claim downgraded to a partial exposure artifact. \\
Timing / leakage & Pooled AUC may reuse target-period information. & Rolling-origin, strict holdout, leakage audit. & PR-AUC, AUC, precision@$k$, recall@$k$. & Retains out-of-time discrimination. & Retrospective ranking only. \\
Case dominance & One large case could drive the result. & Leave-one-case-out and leave-largest-case-out. & Distribution of PR-AUC, AUC, recall. & Stable across cases. & Case-composition driven. \\
Legal-scope asymmetry & Must not be a generic membership detector. & Direct-defendant classification. & AUC and PR-AUC near random. & Not a membership detector (loser-side scope confirmed). & Scope claim unclear. \\
Economic profile & Target should have economic content. & Firm-level and bid-level profile comparisons. & Standardized differences, bid-gap moments. & Cobidders differ on economic dimensions. & Target has less economic content. \\
Bid-layer complementarity & Award layer must add non-redundant information. & Award, bid, and combined models on the same target. & PR-AUC, AUC, precision, recall. & Award adds non-redundant signal. & Award layer redundant. \\
Gatekeeping & Sequencing value must hold at affordable budgets. & Cost-recall frontier over $K_1$ and cost denominators. & Recall, precision, false positives, cost per true positive. & Concentrates the pool at acceptable recall. & Sequencing value smaller or budget-specific. \\
Strategic response & Published rules invite gaming. & Bunching, refreshed-queue, and adaptation analysis. & Threshold bunching, rank instability. & Ranks and refreshed queues robust. & Prefer continuous ranks and refreshed protocols. \\
\bottomrule
\end{tabular}
\begin{tablenotes}
\scriptsize
\item \textit{Notes:} This is a pre-analysis validation-threat
map, not a results table. It states the tests and the
interpretation if each test passes or weakens; it does not report
outcomes. Predictions evaluate whether the award-layer ranking is
useful for evidence allocation; they do not define legal
membership or proof. The realized results of these tests are
reported in the validation sections and the appendix.
\end{tablenotes}
\end{threeparttable}
\end{adjustbox}
\end{table}

Together, these tests evaluate the screen at the stage for which
it is designed: the allocation of costly forensic attention
before liability evidence is assembled. The BEC cutoff is an
implementation detail; the portable object is the staged
architecture that uses cheap administrative records to discipline
where costly cartel forensics should begin.
